\scanpage{193}
les $\xi_{\ell}$ du (i) dépendent du choix d'un pt \unsure{géom|géom.} de $X$, mais à isom non \unsure{uniques|unique} près (\unsure{définir|défin.|dépendance} \unsure{ens des|ens de|du} \uuline{$\unsure{t_{\alpha}|{\pi}_1}(X)$}] \inserted{[$\pi_1(X)$ opère sur $\xi_\ell$]} sont \unsure{indépendants} de ce choix.

Les $\xi_p^\alpha$ dépendent du choix d'un \deleted{corps} pt \inserted{de $X$} à valeurs dans un corps parfait sans plus, si on veut les $k(x)^{\unsure{\mathrm{prf}|\mathrm{pf}}}$\deleted{\unsure{,}} (clôture parfaite de $k(x)$), et pour $k$ donné, il ne \uline{semble pas} qu'il y ait \deleted{\unsure{aie}} des isom entre les $\xi_{p,a}^{(\alpha)}$ correspondant aux divers pts \unsure{à} valeurs dans $k$ (c'est la ``partie \unsure{continue}'' des $\unsure{H^1|H}$ p.ex. \unsure{qui va|que va} faire des ennuis), s'ils sont localisés en le $\unsure{=|\text{même}}$ $x\in X$, \deleted{\unsure{il est}} \deleted{\unsure{sont}} si $k$ est alg. clos, ils sont conjugués par autom. de $k$ \inserted{(étendus à $\unsure{\mathbb{W}|W}_p(k)$)} les groupes alg. associés sont donc conjugués sous autom de $k$ ---- il est $=$ gentil qu'ils veuillent être l'un tordu de l'autre \ldots\ Le groupe $\pi_1(X)$ n'opère plus sur $\xi_p^\alpha$,
\margin{En revanche, on a comme structure supplémentaire les opérations $F,V$ avec $FV=p$}
par contre le groupe des autom de $k$ sur $k(x)$ \deleted{\unsure{[ou = de $k$|[ou = de $k$,}}
\deleted{\unsure{il existe]}} opère de façon $1/2$ \uline{linéaire} sur $\xi_p^{(\alpha)}$ ---- [\unsure{ou|on} \unsure{une|on} \unsure{donnée de|devient|devine}
\margin{élégante. D'autre part, un autom $\sigma$ de $k$ qui n'est pas $k(x)$-linéaire transforme en général $\xi_{x,p}^{\unsure{\ast|\alpha}}$ en un $\xi_{xp}^{\sigma}$ qui n'est plus isom à $\xi_{x,p}$ ----}

Notations
\diagram{\(\begin{array}{ll}
\text{(i)}&
\left\{\begin{array}{ll}
\xi_{\bar{x},\ell} & \text{foncteur fibre $l$-adique en}\\
 & \text{le pt géom $\bar{x}$}\\
\xi_{x,\unsure{\ell|\beta}} & \text{par abus de langage pour $x\in X$}\\
 & \text{ou un pt de $X$ à val.}\\
 & \text{dans un corps quelc.}\\
 & \text{($\xi_{x,\ell}$ déf. à isom \uline{non unique} près)}
\end{array}\right.\\[1ex]
\text{(ii)}&
\left\{\begin{array}{ll}
\xi_{x,p}&\text{$x$ \deleted{\unsure{est}} pt de $X$ à val. dans un corps}\\
\xi_{x,p}^{0}&
\end{array}\right.
\end{array}\)}
\margin{\uline{On peut mettre} \uline{des exposants $\alpha$} \uline{à volonté}}
\margin{192}
